\chapter{Project Management}
\section{Project Plan}
The scope of the project is to develop a research proposal that designs a context-aware multimodal misinformation detection framework with DANES/GETAE models as baselines.
The proposal will include an introduction of existing misinformation detection approaches, identification of research problem and gap, formulation of research questions,
literature review on context-aware and multimodal framework on misinformation detection, a proposed methodology involving modified DANES/GETAE models,
and determined evaluation metrics on model performance.

This project will not cover the full implementation of DANES/GETAE models proposed in the methodology section, model testing or performance benchmarking.
It also will not cover dataset collection or processing, nor model deployment. Focus of this project is only on developing a research design.

To achieve the project outcome, the team will adopt an Agile-Scrum approach to ensure clear task allocation, accountability and timely development of the research proposal.
Each member of the team will have responsibility to deliver section(s) of proposal.
Project work will be organised into weekly sprints where reguler team meeting will be held for sprint planning, progress update and sprint review to resolve any problem or progress blockers.

Several tools will be used to support collaborative work and project management, such as:
(1) Microsoft Teams for team communications and task distributions,
(2) GitHub for version control and tracking proposal document changes, and
(3) Overleaf for collaborative writing and integration of proposal document.
These tools enable efficient project progress tracking and team communication throughout the project.

\section{Timeline}
Project timeline below is derived from a Work Breakdown Structure (WBS) from the objective of research proposal development.
It divides the project into key phases, including literature review, research design, methodology development, proposal drafting, and finalisation.
The timeline also identifies critical paths, such as the completion of literature review and research design, which are essential for the subsequent development of the proposal.
The project will be conducted over the semester (12 weeks) with the following key milestones:

\begin{table}[H]
\centering
\small
\begin{tabular}{|p{6cm}|c|c c c c c c c c c c c c|}
\hline
\textbf{Milestone} & \textbf{CP} & \textbf{W1} & \textbf{W2} & \textbf{W3} & \textbf{W4} & \textbf{W5} & \textbf{W6} & \textbf{W7} & \textbf{W8} & \textbf{W9} & \textbf{W10} & \textbf{W11} & \textbf{W12} \\
\hline

Topic refinement and preliminary literature review & Y & \cellcolor{blue!20} &  &  &  &  &  &  &  &  &  &  &  \\
\hline

Comprehensive literature review, identification of research gap, and initial formulation of research questions & N &  & \cellcolor{blue!20} & \cellcolor{blue!20} &  &  &  &  &  &  &  &  &  \\
\hline

Finalisation of research questions & Y &  &  &  & \cellcolor{blue!20} &  &  &  &  &  &  &  &  \\
\hline

Design and development of proposed methodology and framework & Y &  &  &  &  & \cellcolor{blue!20} & \cellcolor{blue!20} &  &  &  &  &  &  \\
\hline

Refinement of methodology and definition of evaluation metrics & N &  &  &  &  &  & \cellcolor{blue!20} & \cellcolor{blue!20} &  &  &  &  &  \\
\hline

Drafting of research proposal sections & N &  &  &  &  & \cellcolor{blue!20} & \cellcolor{blue!20} & \cellcolor{blue!20} & \cellcolor{blue!20} & \cellcolor{blue!20} & \cellcolor{blue!20} &  &  \\
\hline

Editing, proofreading, and final submission & Y &  &  &  &  &  & \cellcolor{blue!20} &  &  & \cellcolor{blue!20} & \cellcolor{blue!20} &  &  \\
\hline

Final presentation preparation of the research proposal & N &  &  &  &  &  &  &  &  &  &  & \cellcolor{blue!20} & \cellcolor{blue!20} \\
\hline

\end{tabular}
\caption{Project Timeline (Gantt-style Representation)}
\end{table}

\section{Task Allocation}
Task allocation is important for ensuring accountability and progress of project with each team member responsible for specific deliverables.
The following table outlines the task allocation for each team member based on the project milestones:

\begin{table}[H]
\centering
\small
\begin{tabular}{|p{5.5cm}|c|c|c|c|c|}
\hline
\textbf{Task} & \textbf{Timothy (Leader)} & \textbf{Vaidehi} & \textbf{Laiba} & \textbf{Krati} & \textbf{Lyris} \\
\hline

Topic refinement and preliminary literature review 
& \checkmark & \checkmark & \checkmark & \checkmark & \checkmark \\
\hline

Identification of research gap and formulation of research questions 
&  &  & \checkmark &  &  \\
\hline

In-depth literature review 
&  & \checkmark &  &  &  \\
\hline

Finalisation of research questions 
&  &  & \checkmark &  &  \\
\hline

Development of methodology and framework 
& \checkmark &  &  &  &  \\
\hline

Refinement of methodology and definition of evaluation metrics 
& \checkmark &  &  &  & \checkmark \\
\hline

Project management and coordination 
&  &  &  &  & \checkmark \\
\hline

Drafting of proposal sections 
& \checkmark & \checkmark & \checkmark & \checkmark & \checkmark \\
\hline

Editing, proofreading, and final submission 
& \checkmark &  &  &  & \checkmark \\
\hline

\end{tabular}
\caption{Task Allocation and Responsibilities}
\end{table}
\vspace{0.3cm}
\noindent \textit{Note: All team members contribute to in-depth literature review.}

\section{Risk and Mitigation}
The project adopts a proactive risk management approach by identifying potential risks and assessing their impact
particularly towards research proposal outcome and timeline. To resolve these risks, we are defining mitigation and contingency
strategies to minimise disruption to project progress. Higher ranked risks stated below will be closely monitored and escalated to project leader for resolution if they occur.
Any unresolved risks will be discussed in weekly meetings to determine response and contigency actions.
If there are any risks that occur during the project that have not been identified, it will be added to develop its mitigation, response and contigency plan.

\begin{table}[H]
\centering
\small
\begin{tabularx}{\textwidth}{c|X|c|X|X|X}
\toprule
\textbf{Rank} & \textbf{Risks} & \textbf{Impact} & \textbf{Mitigation} & \textbf{Response} & \textbf{Contingency} \\
\midrule

1 & Unclear research problem or gap leads to poorly defined research questions 
& High 
& Conduct comprehensive literature review on the theme of misinformation detection models 
& Reassess research direction and go through the highlights of conducted literature reviews 
& Clarify objective of research, narrow scope of research and sharpen research questions \\

\hline

2 & Complex methodology design without any benchmark on existing detection models 
& High 
& Align design with proposal scope, referencing / benchmark prior research papers 
& Simplify framework / parameters of models and reduce components 
& Adopt baseline models in prior research paper with modified parameters \\

\hline

3 & Difficulty in selecting datasets for social-context aware multimodal analysis 
& High 
& Identify validated datasets that have been used in prior research papers 
& Evaluate alternative datasets available today that may be compatible with the proposed model 
& Choose alternative datasets aligned with proposed model features and parameters or adjust models to selected datasets \\

\hline

4 & Time constraints lead to incomplete research proposals 
& High 
& Set up weekly milestones and weekly progress tracking 
& Reallocate tasks to different team members with available bandwidth 
& Set priority on core sections and establish a strict timeline (days) \\

\hline

5 & Inconsistent writing quality across members 
& Medium 
& Assign clear section ownership to each member 
& Conduct peer review and standardise writing style 
& Perform final integration and editing by designated member \\

\hline

6 & Uneven workload distribution among team members 
& Medium 
& Clear task allocation and perform weekly progress tracking 
& Reassign tasks during weekly meetings / sprint reviews 
& Reassign tasks during weekly meetings / sprint reviews to adjust future workload \\

\bottomrule
\end{tabularx}
\caption{Risk Management Plan}
\end{table}





\section{DELETE original question}
project management - 1 page (1 mark)

This section should outline how you plan to manage your research project throughout the semester. 

It is essentially
a Project Plan that clearly defines the scope and deliverables, including explicit exclusions. 

It should present
a structured schedule with milestones and a transparent critical path, while assigning responsibility to specific
2individuals to ensure accountability. 

It must also incorporate risk management to anticipate threats and formulate
appropriate responses, protecting project objectives from avoidable failure. 

This section should also clearly state
the methods you will use to manage the project, such as Agile sprints, Scrum ceremonies, Kanban, etc., and
specify the enabling tools (e.g., MS Project, Trello, JIRA, Slack, etc.) that will operationalise these methods.

Briefly explain how these methods and the supporting tools together will facilitate scheduling, progress tracking,
communication, and reporting. Specifically, you must include the following:
\begin{itemize}
\item High-level scope and deliverables of your research project. Clearly state exactly what you will deliver and
explicitly rule out what you will not.
\item Timeline with major milestones derived from the Work Breakdown Structure (WBS). This should reflect
the schedule of your research project and include any critical paths. Note: you don't need to include the
WBS itself.
\item Clear assignment (ownership) of each task to project team member(s), not roles, for accountability and
monitoring.
\item Approach in risk management, identify risks and proposed mitigation and contingency plan, as well as
escalation paths. List the major risks, rank them, and state what you will do before, when, and after they
occur.
\end{itemize}

\section{aaaSection2}
\subsection{aaaSection2Subsection1}
aaa


